\section{Lean formalization and reproducibility}\label{ap:lean}

\subsection{Scope of the certificate}\label{lv:scope}
The supplementary Lean development verifies the global quantitative
orthant comparison with its explicit constants, the separate centroid
penalty, and their consequences for decoding, Procrustes recovery, rank
deficiency, unused energy and mean width.  This extends the earlier
certificate for the ordered channel-parameter crossings, the exact
strict-normal minimizer on sufficiently small Procrustes shells, and the
quantified $O(\rho^6)$ Taylor remainder.  All statements use the actual
Gaussian measures, integrals and geometric distance.  The global
comparison does not assume a witness bound, a product deficit or a
desired CDF inequality: these are proved in the dependency chain,
including for singular covariances.

The global results hold for $m\ge2$, with $\lambda>0$ where a positive
decoding coefficient is required.  This does not relax the separate
hypotheses of the local results: those retain $m\ge3$, the stated strict
curvature gap, and a sufficiently small shell radius.  The certificate
does not cover numerical diagnostics, learning experiments or historical
priority claims.  It does not identify finite-radius minimizers in the
tangential regimes or on a boundary of equal curvatures.  The local
radius and remainder constant are proved to exist; no effective numerical
value is certified.

The final library entry is \path{SimplexFull/VerifiedGlobal.lean}.
Its dependency closure contains 216 source modules, including
52 vendored foundation modules.  A separate audit imports this entry
and prints the axioms of 91 selected public declarations.
These counts describe the supplied artifact, rather than a claim that
every sentence of the manuscript has a formal counterpart.  The formal
and written proofs may use different intermediate lemmas.
The original 94-module local development is retained byte for byte.

\subsection{Objects and normalization}\label{lv:objects}
The Gaussian density and distribution function are concrete definitions:
\path{SimplexFull.phi} is Mathlib's real standard Gaussian density, and
\path{SimplexFull.Phi} is its Lebesgue integral over $(-\infty,c]$.
The definitions of $q,h,\sigma$ therefore have their usual analytic
meaning.  For $m\geq3$, the normalized moments in the crossing theorems are
defined using
\[
 w_{m,\mu}(c)=e^{\mu c}\phi(c)^2\Phi(c)^{m-2},
 \qquad
 Z_m(\mu)=\int_{\mathbb R}w_{m,\mu}(c)\,dc,
\]
and
\[
 a_m(\mu)=\frac{\int h(c)^2w_{m,\mu}(c)\,dc}{Z_m(\mu)},\qquad
 b_m(\mu)=\frac{\int \sigma(c)w_{m,\mu}(c)\,dc}{Z_m(\mu)},\qquad
 H_m(\mu)=\frac{\int h(c)w_{m,\mu}(c)\,dc}{Z_m(\mu)}.
\]
These are respectively \path{SimplexFull.a}, \path{SimplexFull.b}, and
\path{SimplexFull.H}; their normalizer is proved positive.  The parameter
in these integrals is $\mu=\lambda/\sqrt{\alpha}$, where
$\alpha=(m-1)/m$.  Thus a crossing expressed in $\mu$ corresponds to the
physical signal amplitude $\lambda=\sqrt{\alpha}\,\mu$.  The variables
$c$, $\mu$, and $\lambda$ are not interchangeable.  For $\lambda>0$,
\[
 a_m(\mu)=A_\lambda/C_\lambda,\qquad
 b_m(\mu)=B_\lambda/C_\lambda,\qquad
 C_\lambda=\frac{\mu e^{-\mu^2/2}}{m}Z_m(\mu)>0.
\]
The value $\mu=0$ used in the normalized crossing arguments is an endpoint
of the moment family, not a division by the zero-amplitude value of
$C_\lambda$.

The objective \path{SimplexFull.decodingSuccess} is defined by
\[
 \Pc(G;\lambda)=\frac{e^{-\lambda^2/2}}{m}
 \int_{\mathbb R^m}e^{\lambda\max_i x_i}\,dN(0,G)(x).
\]
Here $N(0,G)$ is Mathlib's actual
\path{ProbabilityTheory.multivariateGaussian}, including the singular
simplex covariance.  Gaussian density identities for positive definite
covariances and the common Gaussian shift identity connect the local
analytic calculations to this definition.  The extended development also identifies this expression with the
operational maximum-likelihood decoder for unit signals.  For signals
of unequal energy it instead uses the class likelihood
$\exp(\lambda\langle x_i,z\rangle-\lambda^2\|x_i\|^2/2)$ and the
maximum-likelihood/Bayes-value identity.  The individual energy correction
is retained when proving the unused-energy result.

In the namespace \path{SimplexFull.Geometry},
\path{simplexSuccess m} is exactly the function
$p_m(\lambda)=\Pc(G_\Delta;\lambda)$.  The symbols $p_m'$ and $p_m''$
in the Taylor statement are represented by \path{deriv} and
\path{deriv (deriv (simplexSuccess m))}, rather than by independent
coefficient variables.  The theorem
\path{SimplexFull.contDiff_simplexSuccess} supplies genuine $C^3$
regularity.  Its proof establishes all-real exponential integrability
of the coordinate maximum and applies Mathlib's moment-generating-function
analyticity theorem.  In fact, the formal regularity result is stronger
than the order needed here.

For configurations, \path{Geometry.procrustesSq X V} is defined as the
real infimum
\[
 \inf_{U\in O(d)}\sum_{i=1}^m\|X_i-U V_i\|^2.
\]
The index type consists of linear isometric automorphisms of the whole
Euclidean space.  Accordingly, \path{Geometry.orthoOrbit X N} expresses
membership in the actual $O(d)$ orbit, including reflections.  The
distance is not replaced by a chosen alignment or an auxiliary matrix
norm.  The normal configuration is proved realizable when $d\geq m$,
and its squared Procrustes distance is proved to equal $\rho^2$.

For a finite configuration $X$, the support of its actual convex hull is
defined as the supremum of $\langle g,y\rangle$ over
$y\in\operatorname{conv}(\operatorname{range}X)$.  The development proves
that this supremum is attained and equals the finite maximum of the
scores.  The normalized sphere measure is the normalization of
Mathlib's \path{volume.toSphere}; it is not defined through a Gaussian
ratio.  The polar-coordinate measure-preserving map, the radial Gaussian
density, and positive homogeneity identify twice its support average (equivalently, the average of
$h_K(u)+h_K(-u)$) with the
Gaussian expression $2\E\max_i\langle g,X_i\rangle/\E\|g\|$.
The denominator is always the Gaussian radial mean in the original
ambient dimension $d$, including when an auxiliary lift uses $d+m$
coordinates.

\subsection{Correspondence with the new theorems}\label{lv:map}
Table~\ref{lv:table} lists the principal entry declarations.  Module paths
are relative to \path{SimplexFull}; all declarations shown below are
fully qualified.  The package's proof map and source index give the
remaining dependency information.

\begingroup
\small
\renewcommand{\arraystretch}{1.18}
\begin{longtable}{@{}>{\raggedright\arraybackslash}p{0.25\textwidth}p{0.70\textwidth}@{}}
\caption{Principal mathematical statements and their Lean entry points.}
\label{lv:table}\\
\toprule
Statement & Module and public declaration\\
\midrule
\endfirsthead
\toprule
Statement & Module and public declaration\\
\midrule
\endhead
Theorem~\ref{thm:global}
&
\path{GlobalComparison.lean}:
\path{SimplexFull.GlobalComparison.orthant_sqrt_comparison},
\path{orthant_linear_comparison}, and
\path{quantitative_lowerOrthant} in the same namespace.
These use the exact $K,\kappa$ and all admissible PSD covariances.
\\
Quantitative rectangle and centroid bounds
&
\path{GlobalRectangleBound.lean}:
\path{SimplexFull.GlobalRectangle.symmetricRectangle_gap};
\path{GlobalCentroidConfig.lean}:
\path{SimplexFull.GlobalCentroid.centroid_penalty}.
\\
Theorem~\ref{thm:code} and Corollary~\ref{cor:operational}
&
\path{GlobalApplicationsCore.lean} and
\path{GlobalApplicationsDecoder.lean}:
\path{SimplexFull.GlobalApplication.decodingSuccess_gram_penalty},
\path{decodingSuccess_procrustes_penalty}, and
\path{decodingSuccess_rank_deficiency} in that namespace.
The centroid and operational unused-energy conclusions are
\path{SimplexFull.GlobalTail.decodingSuccess_centroid_penalty} and
\path{SimplexFull.GlobalTail.ball_decoder_unused_energy}.
\\
Theorem~\ref{thm:width}
&
\path{GlobalSphereGaussianWidth.lean}:
\path{SimplexFull.SphereGaussian.ball_mean_width}.
The traditional width-function integral is identified by
\path{GlobalSphereGaussianSymmetry.lean}:
\path{SimplexFull.SphereGaussian.widthIntegral_convexHull_eq_meanWidth}.
The Gaussian representation is
\path{SimplexFull.GlobalApplication.ball_mean_width}.
\\
Theorem~\ref{thm:ordered-switches}
&
\path{GaussianCrossings.lean}:
\path{SimplexFull.gaussian_ordered_crossings}.
The strict exclusion of the intermediate coefficient from the minimum is
\path{SimplexFull.gaussian_middle_coefficient_not_min}.
The complete chart and its positive-scaling form are in
\path{CrossingSignChart.lean}:
\path{SimplexFull.gaussian_ordered_crossings_sign_chart} and
\path{SimplexFull.gaussian_ordered_crossings_sign_chart_scaled}.
\\
Theorem~\ref{thm:triangular-switch}
&
\path{GaussianCrossings.lean}:
\path{SimplexFull.gaussian_triangular_unique_crossing}.
The equivalent equation $\mu H_3(\mu)=1$ is established by
\path{SimplexFull.gaussian_triangular_coefficient_eq_iff}.
\\
Theorem~\ref{thm:normal-shell}
&
\path{UniformNormalShell.lean}:
\path{SimplexFull.Geometry.gaussian_uniform_exact_normal_shell_m4}
and
\path{SimplexFull.Geometry.gaussian_uniform_exact_normal_shell_m3}.
The corresponding threshold versions are
\path{SimplexFull.Geometry.gaussian_normal_threshold_uniform_shell_m4}
and
\path{SimplexFull.Geometry.gaussian_normal_threshold_uniform_shell_m3}.
\\
Theorem~\ref{thm:normal-taylor}
&
\path{GaussianTaylor.lean}:
\path{SimplexFull.Geometry.normalTaylor_bound},
\path{SimplexFull.Geometry.normalTaylor_isBigO},
\path{SimplexFull.Geometry.gaussian_uniform_taylor_shell_m4},
and
\path{SimplexFull.Geometry.gaussian_uniform_taylor_shell_m3}.
\\
\bottomrule
\end{longtable}
\endgroup

For $m\geq4$, the crossing declaration asserts
$0<t_1<t_0<t_{01}$ and uniqueness on $[0,\infty)$ of the equations
\[
 a_m(t_1)+(m-2)b_m(t_1)=m-1,\qquad
 a_m(t_0)=m/2,\qquad
 a_m(t_{01})=m b_m(t_{01}).
\]
The declaration \path{gaussian_ordered_crossings_sign_chart} packages
these roots and their uniqueness together with all four open-interval
comparisons and all three endpoint comparisons for $L_0,L_1,L_\perp$.
Its \path{gaussian_ordered_crossings_sign_chart_scaled} companion gives the
same chart after multiplication by any positive common factor, with the
roots independent of that factor.  For every $\lambda>0$ one may therefore
use $C_\lambda>0$ and $\Lambda_j=C_\lambda L_j$ to obtain the physical
coefficient chart in Theorem~\ref{thm:ordered-switches}.  No additional
analytic hypothesis is required.  The original crossing declaration
instantiates the abstract comparison lemmas with the actual Gaussian moments.
The proofs include the strict benchmark $b_4(0)>1/2$,
monotonicity, and the integral balance identity
$2a_m+(m-1)b_m=m+\mu H_m$.
The case $m=3$ is handled separately: its crossing is
$a_3+b_3=2$, and no nonzero $\mathcal H_0$ sector is introduced.

The geometric conclusions are packaged as
\path{SimplexFull.Geometry.UniformExactNormalShell} and
\path{SimplexFull.Geometry.UniformTaylorNormalShell}.
For fixed $m$ and $\lambda$ in the strict-normal regime, the latter has
the quantifier order
\[
 \exists C>0\ \exists\rho_0>0\
 \forall d\geq m\ \forall V\text{ with Gram }G_\Delta\
 \exists z\
 \forall\,0<\rho<\rho_0\
 \forall X\text{ on the feasible shell}.
\]
The normal vector $z$ is unit and orthogonal to every $V_i$.
In particular, both $C$ and $\rho_0$ are independent of the ambient
dimension, reference embedding, and competitor.  The chosen $z$ is
independent of $\rho$.  The final statements contain feasibility of
$N_i=\kappa V_i+\sqrt{1-\kappa^2}\,z$, the exact optimal value, a comparison
with every feasible $X$, and equality if and only if
\path{Geometry.orthoOrbit X N}.  The strict-normal hypothesis is
$a_m(\mu)>m/2$ for $m\geq4$, and $a_3(\mu)+b_3(\mu)>2$ for $m=3$.

The Taylor remainder is the actual function
\[
 \mathcal R_{m,\lambda}(\rho)
 =p_m(\lambda)-p_m\!\left(\lambda\left(1-\frac{\rho^2}{2m}\right)\right)
 -\left[
 \frac{\lambda p_m'(\lambda)}{2m}\rho^2
 -\frac{\lambda^2p_m''(\lambda)}{8m^2}\rho^4
 \right].
\]
The declaration \path{normalTaylor_bound} proves
$|\mathcal R_{m,\lambda}(\rho)|\leq C|\rho|^6$ for $|\rho|<r$,
with $C,r>0$.  The companion \path{normalTaylor_isBigO} states the
corresponding two-sided big-$O$ assertion at zero.  The generic proof
uses a third-order Taylor expansion, obtains an $O(h^3)$ quadratic
remainder, and substitutes $h=-\lambda\rho^2/(2m)$.  Thus it proves
$O(\rho^6)$, rather than only the $o(\rho^4)$ conclusion available
from $C^2$ regularity.

The stronger regularity declaration
\path{SimplexFull.contDiff_simplexSuccess_top} also verifies the
analyticity assertion.  In the pinned Mathlib version the order type of
\path{ContDiff} is $\operatorname{WithTop}(\mathbb N_\infty)$:
its outer $\top$ is the analytic order $\omega$, distinct from the
smooth order $\infty$.  Applying \path{ContDiffAt.analyticAt} therefore
gives \path{AnalyticAt} at every real argument.  A supplementary
kernel-checked statement spells out this consequence for the actual
multivariate Gaussian probability integral.

The declaration
\path{SimplexFull.Geometry.normalDeficit_odd_iteratedDeriv_zero}
in \path{TaylorParity.lean} proves that every odd derivative of the exact
deficit vanishes at zero.  Finally,
\path{SimplexFull.Geometry.simplex_quadratic_coefficient_eq_normal_curvature}
in \path{TaylorCoefficient.lean} identifies the leading coefficient with
the actual normal curvature:
\[
 \frac{\lambda p_m'(\lambda)}{2m}
 =\frac{m-1}{2}\,C_\lambda.
\]
This identity includes the conversion between $\lambda$ and $\mu$.


\subsection{Global proof chain and assumptions}\label{lv:global}
The scalar function $\Psi$ is defined from the actual inverse of
$h(s)=s+q(s)$ on the positive axis.  Its monotonicity, inverse range,
strict concavity, and derivatives are proved.  The curvature constant
$M_0$ is the attained maximum of the actual continuous curvature
function on the manuscript's compact interval.  In particular,
the lower Taylor bound used in the centroid argument is not supplied
as an external smoothness or remainder assumption.

The quantitative rectangle proof starts from the true bivariate
conditional Gaussian law.  Differentiation under the compact integral
and the boundary derivative give the Plackett formula.  Integration
on $[0,1/2]$, monotonicity and evenness yield the stated
$\ell^2e^{-4T^2/3}/(4\pi)$ coefficient.  At correlation $+1$ or $-1$
the residual Gaussian is a Dirac measure and the rectangle mass is
the smaller marginal slab mass; no singular density is substituted.
The actual Gaussian slab peeling argument and a maximizing
off-diagonal pair then give the factor $m(m-1)$ and the exponent $m-2$.

For Theorem~\ref{thm:global}, existence of actual adaptive witnesses,
their compact coordinate bounds, the sum--difference Gaussian law
and Fubini's theorem lead to the quantitative doubled-product deficit.
The product of the normalized self-convolution factors has Gaussian
integral at least one, including
for PSD singular covariance by continuous regularization.
The fixed-factor rectangle estimate on the central box gives precisely
$K$, with exponent $2m-2$.  The correlation-matrix bound
$\|R-I\|_F^2\le m(m-1)$ gives the manuscript's $\kappa$ by
rationalizing the square root.  The identity covariance is handled
directly.  The centroid proof uses the actual PSD variational value,
the verified lower Taylor bound for $\Psi$, and the exact $\tau,\eta$.

The integration bridge proves
\[
 \E\max Z-\E\max X
 =\int_{\mathbb R}\bigl(F_X(c)-F_Z(c)\bigr)\,dc
\]
from a centered signed-tail kernel and Fubini, with integrability
proved.  Applied to the exponential of the maximum it retains the
exact window factor $e^\mu-1$.  Gaussian normalization at
$R=\alpha G+J/m$ supplies the factors $\alpha^2$ and $\sqrt\alpha$.
Consequently the final entries use
\[
 c_g=W\kappa\alpha^2,\quad c_s=W\eta,\quad
 c_P=\left(\frac{\alpha}{c_g}+\frac{1}{mc_s}\right)^{-1},
 \qquad b_g=\kappa\alpha^{3/2},\quad b_s=\eta/\sqrt\alpha.
\]
None of those entries retains a unit-comparison or target-CDF hypothesis.

The geometric inequality is proved for vectors of arbitrary lengths
and for the actual infimum over $O(d)$; a polar-alignment witness also
proves attainment.  The rank corollary assumes the actual matrix rank
bound, with its bound independent of any realization dimension.
For ball configurations the formal proof uses mutually orthogonal
extra coordinates.  Adding the resulting independent centered Gaussian
noise cannot decrease the expected maximum; the corresponding likelihood
argument proves the operational decoding comparison for nonunit signals.
This orthogonal lift is an alternative to the shared-coordinate lift
in the written unused-energy proof.

The foundation sources from the weak-simplex and simplex-uniqueness
developments are supplied with exact upstream commits, per-file
provenance and their original licenses and attributions in
\path{VENDOR_PROVENANCE.json} and \path{licenses/}.  Their imported
closure is included in the same fresh trust-level-zero build as the
new modules, rather than accepted through inherited project binaries.

\subsection{Build and axiom audit}\label{lv:certificate}
The released archive is
\path{concrete-simplex-lean-global-20261106.zip}.
Its SHA-256 digest is
\begin{quote}\small
\path{4e8e24831e8c56ab07523aed22c89c4ba6854726accccff2aa32c0587133896b}.
\end{quote}
The toolchain is Lean 4.31.0, pinned by
\path{leanprover/lean4:v4.31.0}; the Mathlib revision is
\begin{quote}\small
\path{fabf563a7c95a166b8d7b6efca11c8b4dc9d911f}.
\end{quote}
The successful recorded run is \path{20261106-180522-470580}.
Its \path{evidence/verification.json} records all 216 source-module
commands and exit codes, source and compiled-file hashes, the successful
combined import, and the final axiom audit.  All source hashes remained
unchanged through verification.  The manuscript bundle also supplies
the independent Grok Build and agent cross-reviews.

The verifier copies the three project-source namespaces into a fresh
snapshot and creates an initially empty output directory.  It compiles
their explicit dependency closure in dependency order.  All inherited
project build paths, including paths containing only a bare root module,
are excluded from \path{LEAN_PATH}; the pinned Mathlib cache is supplied
by Lake.  Every project source is compiled with Lean's trust level zero:
\begin{verbatim}
lean -j 2 -t 0 -R <snapshot/src> \
  -o <fresh/build/Module.olean> \
  <snapshot/src/Module.lean>
\end{verbatim}
The final audit imports \path{SimplexFull.VerifiedGlobal} and invokes
\path{#print axioms} for 91 selected declarations.
Their dependency reports contain only
\path{propext}, \path{Classical.choice}, and \path{Quot.sound}.
Acceptance requires exactly one recognized report per requested
declaration.  Missing, duplicate, unexpected or nonstandard reports
cause verification to fail.

After removing comments and strings, source scanning rejects
\path{sorry}, \path{admit}, project \path{axiom} declarations,
\path{unsafe}, \path{native_decide} and \path{debug.skipKernelTC}.
Import parsing rejects unknown project roots and unsupported header
layouts instead of silently omitting them.  Positive and negative
regression tests cover the parser, binary-path isolation and axiom
reports.  These checks accompany, rather than replace, kernel checking.

\subsection{Reproducing the verification}\label{lv:reproduce}
After extracting the archive, install the pinned Lean toolchain with
\path{elan} and make Python available.  From the package directory run
\begin{verbatim}
lake update
lake exe cache get
lake env python reproduce.py
\end{verbatim}
The Mathlib revision is fixed in \path{lakefile.toml}.
The verifier uses \path{SimplexFull.VerifiedGlobal} and
\path{AUDIT_THEOREMS_GLOBAL.json}, creates a fresh timestamped build,
and returns a nonzero status on any source, compilation or axiom-audit
failure.  The package supplies \path{PROOF_MAP_GLOBAL.md},
\path{SOURCE_INDEX.json}, \path{SHA256SUMS.json},
the individual compiler logs, upstream provenance and independent
cross-reviews.  The source, compiler output and
\path{evidence/verification.json} are the reproducibility evidence.
